\chapter{Einf�hrung}
Obwohl schon viel gearbeitet und geforscht wurde auf dem Gebiet der Fluidsimulation, steckt sie immer noch in den Kinderschuhen.
Insbesondere gibt es noch keine interaktiven Verfahren, die auf gro�en Gebieten Wasser im Raum, also nicht als H"ohenfeld realisieren. Anwendungen hierf�r w�ren etwa Wasserf�lle, Getr�nke oder Wasserwerfer in Computerspielen, aber auch zum Beispiel Blut in Operationssimulatoren.

Den Ansto� zu einem um Gr"o�enordnungen umfangreicheren Simulationsgebiet brachte Herr Prof. Dr. Westermann, als er vorschlug, Zellinformationen nur noch beim Partikel zu speichern und nach diesem Zellindex die Partikel zu sortieren. Durch diese implizite Anordnung der Zellen wird erreicht, dass der Speicherbedarf unabh"angig von der Gr��e des simulierten Gebiets wird und meine Software keine 15MB ben�tigt, um 3000 Partikel in einem Gebiet bestehend aus $(100.000)^3$ Zellen zu simulieren. M�chte man 500.000 Partikel verwenden, so stehen immernoch $(16.000)^3$ Zellen zur Verf�gung.

Laut Nick Fosters \& Ronald Fedkiws {\em "`Practical animation of liquids"'} \cite{FOFE} wurde z.B. f�r die Simulation der Schlammdusche von Dreamworks "`Shrek"' lediglich ein Gitter von 150 x 200 x 150 Zellen verwendet.
Um auf meiner Datenstruktur Physik berechnen zu k"onnen und vor allem um die zugeh"orige Oberfl"ache zu konstruieren, sind allerdings sehr spezialisierte Ans"atze gefragt.
So orientiert sich diese Arbeit im Groben an \cite{ETHZ} von Matthias M"uller, David Charypar \& Markus Gross, passt deren Ans"atze aber auf die erw"ahnte Datenstruktur an.
Au"serdem wird die auf ATI-Grafikkarten vorhandene TRUFORM$^\texttt{TM}$-Erweiterung genutzt, um auf einem relativ groben Gitter mittels des Marching Cube Algorithmus erzeugte, durch ebene Teilfl"achen charakterisierte, also eckige,
Oberfl"achen rund darzustellen.

\section{Zweck und Aufbau dieser Diplomarbeit}
W�hrend urspr�nglich geplant war, aus den Ergebnissen aus \cite{ETHZ} und \cite{UTAH} eine interaktive Simulation f�r Fl�ssigkeiten mit verschiedenen Eigenschaften zu entwickeln, entstanden aus der "`entdeckten"' Methode zur Kollisionsdetektion ganz neue Anforderungen und die Physik trat zun�chst in den Hintergrund.
Au�erdem hat diese Arbeit m"oglicherweise auch Relevanz f�r astronomische Simulationen, wof�r die Methoden der Smoothed Particle Hydrodynamics (SPH) urspr�nglich entwickelt wurden, da hier Verfahren aufgezeigt werden, die den Simulationsbereich herk�mmlicher Anwendungen um Gr��enordnungen �bersteigen.

So werden in Kapitel \ref{Physik} die physikalischen Grundlagen erarbeitet, sowie die Smoothed Particle Hydrodynamic als Methode zur numerischen L�sung von Flussproblemen vorgestellt. In Kapitel \ref{KOL} wird auf die f�r SPH ben�tigte Kollisionsdetektion eingegangen und ein neues Verfahren eingef�hrt. In Kapitel \ref{vis} werden Methoden zur schnellen Oberfl�chenerzeugung entwickelt. Kapitel \ref{anwendung} stellt ein kleines Demonstrationsprogramm vor, welches sich auch auf der beigef�gten CD befindet. In Kapitel \ref{resultate} sind einige Screenshots, so wie Angaben zur Performance zu finden. Kapitel \ref{ausblick} gibt schlie�lich einen Ausblick auf Anwendungsm�glichkeiten und Entwicklungspotential.